\documentclass[aspectratio=169,11pt]{beamer}
\usepackage[utf8]{inputenc}
\usepackage[T1]{fontenc}
\usepackage{lmodern}
\usepackage[spanish,es-noshorthands]{babel}
\usepackage{booktabs,graphicx,amsmath,xcolor,tikz}
\usetikzlibrary{positioning,arrows.meta}
\usepackage{array,amssymb}
\newcolumntype{L}[1]{>{\raggedright\arraybackslash}p{#1}}
\tikzset{box/.style={draw, rounded corners, fill=#1!12, minimum width=28mm, minimum height=9mm, align=center}}
\usetheme{default}
\usecolortheme{seahorse}
\setbeamertemplate{navigation symbols}{}
\setbeamertemplate{footline}[frame number]
\definecolor{ares}{HTML}{B3261E}
\setbeamercolor{frametitle}{fg=ares}
\setbeamercolor{title}{fg=ares}
\setbeamercolor{structure}{fg=ares}

\InputIfFileExists{generated/numbers.tex}{}{}
\newcommand{\val}[1]{\ifcsname #1\endcsname\csname #1\endcsname\else\textbf{\textcolor{red}{??}}\fi}
\newcommand{\inputgen}[1]{\IfFileExists{generated/#1}{\input{generated/#1}}{\textbf{\textcolor{red}{[\detokenize{#1} pendiente]}}}}
\newcommand{\figgen}[2]{\IfFileExists{figs/#1}{\includegraphics[width=#2]{figs/#1}}{\fbox{\parbox{#2}{\centering\textcolor{red}{[figura \detokenize{#1} pendiente]}}}}}

\title{ARES-U}
\subtitle{Evolución generacional de un agente neuronal que transfiere patrones entre cinco generaciones de Pokémon}
\author{Cristóbal Moraga --- Rol 202130008-2}
\institute{IPD434 Seminario de Soft Computing --- UTFSM}
\date{Diciembre 2026}

\begin{document}

\begin{frame}
\titlepage
\centering\small Código: \url{github.com/crismoraga/ARES-PKMN} (rama \texttt{ipd434})
\end{frame}

\begin{frame}{El problema y la práctica actual}
\begin{columns}[T]
\column{0.55\textwidth}
\begin{itemize}
\item \textbf{RL} (PPO sobre \emph{Red}): un solo juego, recompensas y scripts específicos.
\item \textbf{LLM + harness} (Gemini, Claude): terminan juegos con RAM a texto y \emph{pathfinders}. \textbf{No aprenden} entre corridas.
\item \textbf{Bots de combate}: una generación, un formato.
\item En 72 trabajos revisados, \textbf{nadie entrena en un juego y evalúa en otro}.
\end{itemize}
\column{0.42\textwidth}
\begin{block}{Propuesta}
Un solo agente que \textbf{reconoce los patrones comunes} de la saga y \textbf{evoluciona} de generación en generación: propuestas que compiten, se fusionan y engendran la siguiente generación.
\end{block}
\end{columns}
\end{frame}

\begin{frame}{Vicios del estado del arte (17 + 6 propios catalogados)}
\small
\begin{tabular}{p{0.30\textwidth}p{0.58\textwidth}}
\toprule
Vicio & Evidencia \\
\midrule
Harness / RAM privilegiada & minimapas, \emph{pathfinders}, RAM a texto (Gemini, Claude) \\
``Sin scripts'' que sí los usa & Jev: lee RAM, A*, hitos a mano, recarga checkpoints (su README) \\
Scripts y trampas & dinero infinito, HMs automáticos (PokeRL ``beat Red'') \\
Cero generalización & todo en un ROM; modelos por tramo \\
N=1, métricas incomparables & horas vs pasos vs turnos; sin semillas \\
\midrule
\textbf{ARES-PKMN (antes)} & solo RAM Gen 3, 101 hitos a mano, red con peso $0.05$ \\
\bottomrule
\end{tabular}
\end{frame}

\begin{frame}{Patrones que persisten, dificultad que crece}
\small
\begin{tabular}{llp{0.62\textwidth}}
\toprule
Gen & Juego (consola) & Se mantiene / se agrega \\
\midrule
1 & Rojo (GB) & cuadrícula, menús, tabla de tipos, gimnasios \\
2 & Cristal (GBC) & + color, reloj día/noche, nuevos tipos \\
3 & Esmeralda (GBA) & + habilidades, combates dobles, mapas más grandes \\
4 & Platino (NDS) & + división física/especial, pantalla táctil, dos pantallas \\
5 & Negro (NDS) & + perspectiva 3D, cámara dinámica \\
\bottomrule
\end{tabular}

\vspace{0.3cm}
Pregunta: \textbf{¿puede una red aprender lo que se mantiene, adaptarse a lo que se agrega y mejorar de generación en generación?}
\end{frame}

\begin{frame}{Un agente, ocho juegos, cuatro consolas}
\begin{columns}[T]
\column{0.5\textwidth}
\figgen{e3_montage.png}{\textwidth}
\column{0.48\textwidth}
\small
\begin{itemize}
\item mGBA multinúcleo (GB/GBC/GBA) y DeSmuME (NDS).
\item Misma red: píxeles $84\times84$, 9 botones, sin táctil.
\item RAM \textbf{solo para medir}; adaptadores descubiertos caminando.
\item Métrica universal (pantallas únicas) validada: $\rho=\val{EthreeRho}$ con celdas RAM.
\end{itemize}
\end{columns}
\end{frame}

\begin{frame}{El ciclo generacional: proponer, combatir, fusionar, heredar}
\begin{columns}[T]
\column{0.46\textwidth}
\centering
\begin{tikzpicture}[node distance=7mm and 9mm, every node/.style={font=\small}, >=Stealth]
\node[box=red] (p) {\textbf{Proponer}\\$K$ hijos mutados};
\node[box=orange, right=of p] (c) {\textbf{Combatir}\\torneo / duelos};
\node[box=green!60!black, below=of c] (f) {\textbf{Fusionar}\\destilación / sopa};
\node[box=blue, below=of p] (h) {\textbf{Heredar}\\padre de G$_{g+1}$};
\draw[->, thick] (p) -- (c); \draw[->, thick] (c) -- (f); \draw[->, thick] (f) -- (h); \draw[->, thick] (h) -- (p);
\end{tikzpicture}
\column{0.52\textwidth}
\small
\begin{itemize}
\item \textbf{Proponer}: 4 hijos del padre; élite + mutaciones de hiperparámetros \emph{(evolutiva)}.
\item \textbf{Combatir}: todos contra todos en el Stadium, rating Elo; rivales congelados en la liga \emph{(adversarial)}.
\item \textbf{Fusionar}: mismo padre $\Rightarrow$ sopa de pesos; distinto origen $\Rightarrow$ destilación del top-3 \emph{(transferencia)}.
\item \textbf{Heredar}: la propuesta única engendra la generación siguiente.
\item Solo validación para elegir; el test nunca se mira.
\end{itemize}
\end{columns}
\end{frame}

\begin{frame}{IPD434 en un solo sistema}
\small\centering
\begin{tabular}{p{0.38\textwidth}p{0.52\textwidth}}
\toprule
Unidad & Componente de ARES-U \\
\midrule
RL: MDP, DQN, PPO (05-1) & agentes de combate y exploración; DQN de línea base \\
Transformers y transferencia (05-2) & Transformer de entidades; ajuste fino entre juegos \\
AE / adversarial (05-3) & encoder AE; \emph{self-play} y liga de rivales \\
CNN avanzadas (05-4) & encoder convolucional transferido entre consolas \\
Computación evolutiva (03, 05) & torneo, mutación, elitismo, fusión de la población \\
\bottomrule
\end{tabular}
\end{frame}

\begin{frame}{¿Cómo se mide? ¿Cómo se sabe si está bien?}
\small
\begin{columns}[T]
\column{0.5\textwidth}
\textbf{Combate}
\begin{itemize}
\item Puntaje $=$ (victorias $+$ empates/2) / combates, frente a Greedy.
\item Combates de \textbf{test} nunca usados para elegir nada.
\item \emph{Está bien} si gana a la referencia con $p_{\mathrm{Holm}}<0.05$ (test pareado: mismos equipos y semillas).
\end{itemize}
\textbf{Exploración}
\begin{itemize}
\item Pantallas únicas por episodio $\to$ normalizado: \textbf{0 = azar, 1 = práctica actual}.
\item \emph{Está bien} si el IC95\% queda $>0$ al llegar (arranque) y $>1$ tras entrenar.
\end{itemize}
\column{0.47\textwidth}
\begin{block}{La medición también se mide}
\begin{itemize}
\item Métrica vs.\ RAM: $\rho=\val{EthreeRho}$.
\item Selección: correlación validación--test.
\item Hallazgos: \textbf{se repiten en datos nuevos}.
\item Lugares: ids de mapa en RAM que el agente nunca ve.
\end{itemize}
\end{block}
Criterios fijados en el código \textbf{antes} de correr (commits \texttt{53d61ec}, \texttt{392d951}, \texttt{3e6e549}, \texttt{3fca0a4}).
\end{columns}
\end{frame}

\begin{frame}{Tablero de veredictos (1/3)}
\centering\footnotesize\renewcommand{\arraystretch}{1.1}
\inputgen{tab_score_a.tex}
\end{frame}

\begin{frame}{Tablero de veredictos (2/3)}
\centering\footnotesize\renewcommand{\arraystretch}{1.1}
\inputgen{tab_score_b.tex}
\end{frame}

\begin{frame}{Tablero de veredictos (3/3)}
\centering\footnotesize\renewcommand{\arraystretch}{1.1}
\inputgen{tab_score_c.tex}
\end{frame}

\begin{frame}{E1: patrones de combate entre generaciones}
\begin{columns}[T]
\column{0.45\textwidth}
\small
\begin{itemize}
\item Tokens de atributos, \textbf{sin IDs} de especie ni movimiento.
\item Transformer pre-norm, atención de puntero sobre acciones legales.
\item Entrena con Gen 1; evalúa en Gen 2/3 \textbf{nunca vistas}.
\item 5 semillas; tests pareados; hipótesis pre-registradas.
\item Supera a Greedy en Gen 3: \textbf{\val{EonePpoGreedyGenThree}} (H1, $p_{\mathrm{Holm}}=\val{EoneHOneGenThreeP}$).
\end{itemize}
\column{0.53\textwidth}
\scriptsize\inputgen{tab_hyp.tex}
\end{columns}
\end{frame}

\begin{frame}{E3: lo aprendido en GBA sirve en GB, GBC y NDS}
\centering\scriptsize
\inputgen{tab_e3.tex}

\vspace{0.2cm}\small
Zero-shot $>$ azar en \textbf{\val{EthreeZeroShotWins}/\val{EthreeGames}} juegos nuevos. Heredar solo el encoder $\approx$ desde cero: \textbf{se transfiere conducta}, no solo percepción.
\end{frame}

\begin{frame}{E4 + E5: el agente evoluciona de Gen 1 a Gen 5}
\centering
\figgen{e4_curriculum.pdf}{0.98\textwidth}

\small Arranque $>0$ en cada juego nuevo (IC95\% $>0$ en 3/3). Tras entrenar supera a la práctica actual con IC95\% $>1$ en Esmeralda y Platino, no en Cristal. Costo: olvido de juegos previos (Rojo \val{EfourFinetuneRojoFinal} al final).
\end{frame}

\begin{frame}{E5-combate: las propuestas pelean, se fusionan y heredan}
\centering
\figgen{e5_battle.pdf}{0.98\textwidth}

\small G0: 10 propuestas de E1 (PPO y BC$\to$PPO) $\to$ torneo $\to$ destilación. G1/G2: 4 hijos en Gen 2 / Gen 3 $\to$ torneo $\to$ sopa.
\end{frame}

\begin{frame}{E5-combate: lo que la medición confirmó y lo que rechazó}
\small
\begin{itemize}
\item \textbf{Se adapta a cada generación y no olvida Gen 1}: Gen 2 \val{EfivebUZeroGenTwo} $\to$ \val{EfivebUOneGenTwo} $\to$ \val{EfivebUTwoGenTwo}.
\item v1: la mejora en Gen 1 (C6, $p_{\mathrm{Holm}}=\val{EfivebCSixP}$) \textbf{no se replicó} en 2000 combates nuevos.
\item Diagnóstico con sus propios datos: torneo de 200 combates por par $\Rightarrow$ error $\pm0.035$; las propuestas difieren $<0.015$ $\Rightarrow$ \textbf{se elegía al azar}.
\item v2 (corrige): nichos, $10^4$ combates pareados, fusión pareada. Resultado: los hijos difieren $\pm0.005$ (error $0.003$) $\Rightarrow$ \textbf{no hay diferencias reales que seleccionar}.
\item Conclusión: la población aporta solo con diversidad real (escuelas distintas + alineación de permutaciones), no con mutaciones de un mismo padre.
\end{itemize}
\end{frame}

\begin{frame}{E6: de pantallas a lugares (sin leer la RAM)}
\small
\begin{columns}[T]
\column{0.52\textwidth}
\begin{itemize}
\item Vicio detectado por el tablero: el agente ve \textbf{más pantallas} pero \textbf{no más lugares} (mapas en RAM).
\item Patrón común a toda la saga: al cambiar de lugar la pantalla \textbf{funde a negro o blanco}.
\item Bono de \emph{novedad de lugar}: fundido $\to$ escena nueva en el episodio. Solo píxeles.
\item Detector validado contra la RAM: detecta el 100\% de los mapas nuevos en Rojo, Cristal y Esmeralda.
\end{itemize}
\column{0.45\textwidth}
\begin{block}{Criterio (fijado antes)}
Más mapas distintos que desde cero, IC95\% $>0$ en $\geq2$ de 3 juegos, sin perder más del 10\% de pantallas.
\end{block}
\textbf{Resultado: no se cumple} (0/3).
\begin{itemize}
\item Coef.\ 1: la señal era ${\sim}1\%$ de la recompensa $\Rightarrow$ sin efecto.
\item Coef.\ 50: \textbf{\emph{reward hacking}} medido: START+A de 3\% a 39\% (Esmeralda); abre menús, que también funden.
\item El tablero lo detectó solo.
\end{itemize}
\end{columns}
\end{frame}

\begin{frame}{Análisis crítico y limitaciones}
\begin{itemize}
\item \textbf{Explorar no es avanzar en la historia}: nada de esto equivale a terminar un juego, y no se afirma.
\item E1/E5-combate cambian especies, no la mecánica de Gen 4+; combates en simulador.
\item 3 semillas en E3--E5; presupuesto NDS corto (DeSmuME ${\sim}150$ cuadros/s).
\item Inicio de Platino con una elección manual en el tutorial (documentada).
\item Con el mismo cómputo la población no rinde más: los hijos de un padre no difieren lo suficiente.
\item Más pantallas no es más lugares; la señal de lugar por fundido fue explotada.
\end{itemize}
\end{frame}

\begin{frame}{Conclusiones y trabajo futuro}
\begin{itemize}
\item La saga comparte patrones que una red aprende y reutiliza: en combate (atributos) y en exploración (conducta).
\item El mismo agente corre en \textbf{8 juegos de 5 generaciones}; en 2 de 3 juegos medidos supera a la práctica actual de entrenar cada juego desde cero.
\item \textbf{Tablero de veredictos}: cada afirmación con criterio previo, IC y replicación; también rechaza (población, replicación, \emph{reward hacking}).
\item El ciclo \textbf{proponer--combatir--fusionar--heredar} hace visible la evolución y une transferencia, aprendizaje adversarial y computación evolutiva.
\item \textbf{Siguiente}: repaso contra el olvido, fusión entre escuelas con alineación de permutaciones, subobjetivos de historia, combate en el cartucho.
\end{itemize}
\vspace{0.3cm}
\centering\large ¿Preguntas?
\end{frame}

\end{document}
